\section{Official values and the 2011 flood: a cross-section}\label{sec:cross}

This section asks a descriptive question: about ten years after the 2011 flood (the 2023--26 values rest on evidence from about 2021--22), how do the official values vary with the share of a cadastral sheet that was inundated? The answer motivates the design of Section~\ref{sec:design} and supplies its power inputs. It is not an estimate of capitalisation, and I give the reasons at the end.

The sample is the core Chao Phraya basin: seventeen provinces from Bangkok north to Phitsanulok and Kamphaeng Phet, restricted to decoded sheets at the two scales that carry most parcels, 1:4,000 (rural) and 1:1,000 (urban). It contains \pilotSheets{} sheets with 2011 flood data and \pilotParcels{} parcels; a further \pilotSheetsNoFlood{} sheets have no flood pixel and are excluded. Appendix Figure~\ref{fig:maps} maps the 2011 flood share and the mean log official value by sheet. The specification compares sheets within a cell defined by province, 1:50,000 map block and scale, so that each flooded sheet is compared with dry sheets a few kilometres away at the same cadastral scale:
\begin{equation}
\bar y_{j} = \sum_{b} \beta_b \,\mathbf{1}[f_j \in b] + \mu_{c(j)} + x_j'\theta + \varepsilon_j,
\label{eq:cross}
\end{equation}
where $\bar y_j$ is the parcel-weighted mean of log value per square wah on sheet $j$, $f_j$ the 2011 flood share, $b$ one of five bins above the reference bin of less than 1\%, $\mu_c$ a cell effect, and $x_j$ optionally the sheet's 2021 built-up, cropland, tree and water shares. Sheets are weighted equally, as in the registered design, and standard errors are clustered by 1:50,000 block.

\begin{table}[t]
\centering\footnotesize
\caption{Official land value by 2011 flood-share bin}\label{tab:binned}
% Generated by 04_analysis/src/paper_numbers.py. Do not edit.
\begin{adjustbox}{max width=\textwidth}
\begin{tabular}{lcccc}
\toprule
 & \multicolumn{2}{c}{Rural sheets (1:4,000)} & \multicolumn{2}{c}{Urban sheets (1:1,000)} \\
\cmidrule(lr){2-3}\cmidrule(lr){4-5}
2011 flood share & (1) & (2) & (3) & (4) \\
\midrule
$<$1\% & 0 (ref.) & 0 (ref.) & 0 (ref.) & 0 (ref.) \\
1--20\% & 0.083 (0.040) & 0.063 (0.023) & $-$0.232 (0.080) & $-$0.019 (0.068) \\
20--40\% & 0.037 (0.043) & 0.076 (0.029) & $-$0.285 (0.039) & $-$0.044 (0.025) \\
40--60\% & 0.048 (0.057) & 0.097 (0.037) & $-$0.381 (0.057) & $-$0.054 (0.038) \\
60--80\% & $-$0.047 (0.056) & 0.100 (0.037) & $-$0.480 (0.066) & $-$0.087 (0.039) \\
80--100\% & $-$0.271 (0.071) & 0.083 (0.043) & $-$0.647 (0.104) & $-$0.083 (0.061) \\
\midrule
Sheets & \multicolumn{2}{c}{12,181} & \multicolumn{2}{c}{28,868} \\
Land-cover controls & no & yes & no & yes \\
\bottomrule
\end{tabular}
\end{adjustbox}

\tabnote{Dependent variable: parcel-weighted mean log official value per \wah\ on the sheet, round 2023--26. Cell effects (province $\times$ 1:50,000 block $\times$ scale) in all columns; columns (2) and (4) add 2021 built-up, cropland, tree and water shares. Sheets weighted equally. Standard errors clustered by 1:50,000 block (\withinClustersRural{} rural, \withinClustersUrban{} urban) in parentheses. Appendix Figure~\ref{fig:binned} plots the coefficients and Appendix Table~\ref{tab:within} gives linear slopes. Source: \texttt{appraisal\_flood\_binned.csv}.}
\end{table}

Table~\ref{tab:binned} gives the main pattern. On urban sheets the raw gradient is monotone and large: sheets flooded 80--100\% in 2011 have mean log values \binTopUrbanNoLc{} (\binTopUrbanNoLcSe) below the reference, about \binTopUrbanPct\% lower, and even sheets flooded 1--20\% are \binLowUrbanNoLc{} (\binLowUrbanNoLcSe) lower. Adding land-cover shares removes most of it: the top-bin coefficient falls to \binTopUrbanLc{} (\binTopUrbanLcSe), and no bin is more than nine log points below the reference; the residual urban gap is about 8\% and imprecise, which is as large as the hazard discounts the literature reports (Section~\ref{sec:status}), not negligible. The reason for the fall is in Appendix Table~\ref{tab:landcover}: urban sheets that stayed dry in 2011 are \lcUrbanBuiltRef\% built-up and \lcUrbanCropRef\% cropland, while those almost wholly flooded are \lcUrbanBuiltTop\% built-up and \lcUrbanCropTop\% cropland. The sheets that the optical satellite saw flooded are, to a large extent, the open land at the edge of towns; the dense centres were under water too but are not in the flooded bins. On rural sheets the raw gradient is weaker and non-monotone: only the top bin is lower, at \binTopRuralNoLc{} (\binTopRuralNoLcSe), and with land-cover controls every bin is positive, between \binLowRuralLc{} and \binTopRuralLc.

Four reasons stop me from reading the cross-section as capitalisation. First, flood-prone land differs from dry land in ways the committees value directly: it is lower, closer to rivers, more often cropland, and farther from roads and town centres. Second, the land-cover control is measured in 2021, ten years after the flood. If the flood itself changed land use, through abandonment, conversion or rebuilding, then land cover is an outcome and a bad control, and the fall from \binTopUrbanPct\% to about 8\% cannot be read as ``mostly land use''; a pre-2011 land-cover product would be needed to separate the two, and none is in the replication files yet. Third, the 2011 flood share is measured with error that is correlated with the outcome: the optical satellite under-detects flooding in built-up areas, which are exactly the high-value sheets, and an urban sheet rests on a median of \supPxOneK{} pixels, so the urban measure is close to binary and noisy. Fourth, there is no pre-flood round in the public data against which to difference, so the cross-section cannot separate what the committees valued before 2011 from what they revised after it.

What the cross-section does establish is useful for the design that follows. The official values carry large flood-correlated variation, so there is something to explain. The within-cell comparison is feasible at scale, the clustering structure is known, and the standard deviation of sheet-level log values within cells sets the scale for the power calculations of Section~\ref{sec:design}. And the land-cover result is a warning that any flood-exposure measure built from satellites must be balanced on land cover and elevation before treated and control sheets are compared.
